% ============================================================
%  Software Requirements Specification
%  GuessTheWord — Progetto Finale JA26
%  Corso di Java Programmazione Avanzata — A.A. 2025/2026
%  Dipartimento DIEM – Università degli Studi di Salerno
% ============================================================

\documentclass[12pt, a4paper]{article}

% ── Pacchetti ──────────────────────────────────────────────
\usepackage[utf8]{inputenc}
\usepackage[T1]{fontenc}
\usepackage[english]{babel}
\usepackage{geometry}
\usepackage{graphicx}
\usepackage{hyperref}

\usepackage{titlesec}
\usepackage{booktabs}
\usepackage[table]{xcolor}
\usepackage{longtable}
\usepackage{array}
\usepackage{enumitem}
\usepackage{fancyhdr}
\usepackage{listings}
\usepackage{tcolorbox}
\usepackage{tabularx}
\usepackage{multirow}
\usepackage{tocbibind}
\usepackage{parskip}



% ── Layout ─────────────────────────────────────────────────
\geometry{
    top=2.5cm, bottom=2.5cm,
    left=2.5cm, right=2.5cm,
    headheight=14pt
}

% ── Colori ─────────────────────────────────────────────────
\definecolor{primaryBlue}{RGB}{30, 80, 160}
\definecolor{lightBlue}{RGB}{220, 235, 255}
\definecolor{darkGray}{RGB}{60, 60, 60}
\definecolor{lightGray}{RGB}{245, 245, 245}
\definecolor{codeGray}{RGB}{240, 240, 240}
\definecolor{accentGreen}{RGB}{0, 120, 60}
\definecolor{accentRed}{RGB}{180, 30, 30}
\definecolor{accentOrange}{RGB}{200, 100, 0}

% ── Titoli ─────────────────────────────────────────────────
\titleformat{\section}
    {\color{primaryBlue}\Large\bfseries}
    {\thesection}{1em}{}[\color{primaryBlue}\titlerule]

\titleformat{\subsection}
    {\color{primaryBlue}\large\bfseries}
    {\thesubsection}{1em}{}

\titleformat{\subsubsection}
    {\color{darkGray}\normalsize\bfseries}
    {\thesubsubsection}{1em}{}

% ── Hyperref ───────────────────────────────────────────────
\hypersetup{
    colorlinks=true,
    linkcolor=primaryBlue,
    urlcolor=primaryBlue,
    citecolor=primaryBlue,
    pdftitle={SRS -- GuessTheWord},
    pdfauthor={Gruppo JA26},
    pdfsubject={Software Requirements Specification}
}

% ── Intestazioni e piè di pagina ───────────────────────────
\pagestyle{fancy}
\fancyhf{}
\fancyhead[L]{\small\color{darkGray} SRS -- GuessTheWord}
\fancyhead[R]{\small\color{darkGray} Corso JA26 -- A.A. 2025/2026}
\fancyfoot[C]{\small\color{darkGray}\thepage}
\renewcommand{\headrulewidth}{0.4pt}
\renewcommand{\footrulewidth}{0.4pt}

% ── Listings ───────────────────────────────────────────────
\lstset{
    backgroundcolor=\color{codeGray},
    basicstyle=\ttfamily\small,
    frame=single,
    rulecolor=\color{darkGray!40},
    breaklines=true,
    captionpos=b,
    showstringspaces=false,
    keywordstyle=\color{primaryBlue}\bfseries,
    commentstyle=\color{accentGreen}\itshape,
    stringstyle=\color{accentRed}
}

% ── Box colorati ───────────────────────────────────────────
\tcbuselibrary{skins, breakable}

\newtcolorbox{reqbox}[2][]{
    colback=lightBlue,
    colframe=primaryBlue,
    title={\textbf{#2}},
    fonttitle=\bfseries\small,
    #1
}

\newtcolorbox{notebox}{
    colback=lightGray,
    colframe=darkGray!60,
    title={\textbf{Nota}},
    fonttitle=\bfseries\small
}

\newtcolorbox{constraintbox}{
    colback=yellow!15,
    colframe=accentOrange,
    title={\textbf{Vincolo}},
    fonttitle=\bfseries\small
}

% ── Comandi utili ──────────────────────────────────────────
\newcommand{\reqid}[1]{\textbf{\texttt{[#1]}}}
\newcommand{\mustshall}{\textbf{DEVE}}
\newcommand{\should}{\textbf{DOVREBBE}}
\newcommand{\may}{\textbf{PUÒ}}
\newcommand{\code}[1]{\texttt{#1}}

% ============================================================
%  DOCUMENTO
% ============================================================
\begin{document}

% ── Copertina ──────────────────────────────────────────────
\begin{titlepage}
    \centering
    \vspace*{1.5cm}

    {\color{primaryBlue}\rule{\linewidth}{1.5pt}}
    \vspace{0.5cm}

    {\Huge\bfseries\color{primaryBlue} Software Requirements Specification}\\[0.4cm]
    {\LARGE\bfseries GuessTheWord}\\[0.2cm]
    {\large Gioco Client–Server a Sfide in Tempo Reale}

    \vspace{0.5cm}
    {\color{primaryBlue}\rule{\linewidth}{1.5pt}}

    \vspace{1.5cm}

    \begin{tcolorbox}[colback=lightBlue, colframe=primaryBlue, width=0.85\linewidth]
        \centering
        \textbf{Corso:} Java Programmazione Avanzata\\[0.2cm]
        \textbf{Codice:} JA26 \quad
        \textbf{Anno Accademico:} 2025/2026\\[0.2cm]
        \textbf{Dipartimento:} DIEM\\
        \textbf{Università degli Studi di Salerno}
    \end{tcolorbox}

    \vspace{1.5cm}

    \begin{tabular}{ll}
        \toprule
        \textbf{Documento}      & Software Requirements Specification (SRS) \\
        \textbf{Versione}       & 1.0 \\
        \textbf{Data}           & Maggio 2025 \\
        \textbf{Standard}       & IEEE 830-1998 (adattato) \\
        \textbf{Stato}          & Definitivo \\
        \bottomrule
    \end{tabular}

    \vfill

    {\small\color{darkGray}
    Questo documento costituisce la specifica formale dei requisiti\\
    per il progetto finale del corso JA26.}
\end{titlepage}

% ── Registro delle revisioni ───────────────────────────────
\newpage
\section*{Registro delle Revisioni}
\addcontentsline{toc}{section}{Registro delle Revisioni}

\begin{longtable}{|c|c|p{6cm}|c|}
    \hline
    \rowcolor{primaryBlue!20}
    \textbf{Versione} & \textbf{Data} & \textbf{Descrizione} & \textbf{Autore} \\
    \hline
    1.0 & Maggio 2025 & Prima emissione. Documento SRS completo derivato dalle specifiche del progetto JA26. & Gruppo \\
    \hline
\end{longtable}

% ── Indice ─────────────────────────────────────────────────
\newpage
\tableofcontents

% ============================================================
\newpage
\section{Introduzione}

\subsection{Scopo del Documento}

Il presente documento costituisce la \textbf{Software Requirements Specification (SRS)} del sistema \textit{GuessTheWord}, sviluppato come progetto finale per il Corso di Java Programmazione Avanzata (JA26), A.A.\ 2025/2026, presso il Dipartimento DIEM dell'Università degli Studi di Salerno.

L'obiettivo di questo documento è fornire una descrizione completa, precisa e non ambigua dei requisiti funzionali, non funzionali e di sistema del software, seguendo le linee guida dello standard IEEE 830-1998. Il documento è destinato al team di sviluppo, al docente valutatore e a chiunque partecipi al processo di revisione del progetto.

\subsection{Scopo del Sistema}

\textit{GuessTheWord} è un \textbf{gioco multiplayer client–server} sviluppato interamente in Java, con interfaccia grafica realizzata in JavaFX. Due utenti si sfidano in tempo reale nel decifrare una o più parole nascoste all'interno di un estratto testuale. Le parole da indovinare vengono cifrate tramite il \textbf{cifrario di Cesare} applicato con uno spostamento scelto casualmente dal server.

Il sistema è composto da tre componenti principali:
\begin{itemize}[leftmargin=*]
    \item \textbf{Server JavaFX}: coordina l'intero sistema, gestisce le connessioni, l'autenticazione, le sfide e la persistenza dei dati.
    \item \textbf{Client JavaFX} (due istanze): forniscono l'interfaccia utente per l'autenticazione, la partecipazione alle sfide e la visualizzazione dello storico.
    \item \textbf{Database SQLite}: garantisce la persistenza di credenziali, sfide e risultati.
\end{itemize}

\subsection{Definizioni, Acronimi e Abbreviazioni}

\begin{longtable}{|p{3.5cm}|p{10cm}|}
    \hline
    \rowcolor{primaryBlue!20}
    \textbf{Termine / Acronimo} & \textbf{Definizione} \\
    \hline
    \endhead
    SRS & Software Requirements Specification \\
    \hline
    JA26 & Codice identificativo del corso \\
    \hline
    JavaFX & Framework Java per la creazione di interfacce grafiche desktop \\
    \hline
    FXML & File XML usato per descrivere layout JavaFX \\
    \hline
    Socket & Meccanismo di comunicazione di rete TCP/IP \\
    \hline
    SQLite & Database relazionale leggero, file-based \\
    \hline
    JDBC & Java Database Connectivity, API standard per accesso a DB in Java \\
    \hline
    Stream API & API Java 8+ per l'elaborazione funzionale di sequenze di dati \\
    \hline
    Service/Task & Classi JavaFX per esecuzione asincrona di operazioni in background \\
    \hline
    Cifrario di Cesare & Tecnica di cifratura per sostituzione di caratteri con uno spostamento fisso \\
    \hline
    TF & Term Frequency: frequenza relativa di una parola in un documento \\
    \hline
    Shift & Valore numerico dello spostamento nel cifrario di Cesare \\
    \hline
    Admin & Utente amministratore con privilegi speciali di gestione del server \\
    \hline
    Sfida & Una partita tra due utenti registrati, coordinata dal server \\
    \hline
    Estratto & Porzione di testo selezionata dal server per una sfida \\
    \hline
    properties & File di configurazione in formato \code{chiave=valore} \\
    \hline
    JAR & Java ARchive, formato eseguibile Java \\
    \hline
    JDK & Java Development Kit \\
    \hline
    RF & Requisito Funzionale \\
    \hline
    RNF & Requisito Non Funzionale \\
    \hline
    RC & Requisito di Consegna \\
    \hline
\end{longtable}

\subsection{Riferimenti}

\begin{itemize}[leftmargin=*]
    \item Specifiche del Progetto Finale — Corso JA26, A.A.\ 2025/2026, Dipartimento DIEM, UNISA.
    \item IEEE Std 830-1998, \textit{Recommended Practice for Software Requirements Specifications}.
    \item Documentazione ufficiale Java 8 — Oracle Corporation.
    \item Documentazione ufficiale JavaFX — OpenJFX Project.
    \item SQLite Documentation — \url{https://www.sqlite.org/docs.html}.
\end{itemize}

\subsection{Panoramica del Documento}

Il documento è strutturato come segue:
\begin{itemize}[leftmargin=*]
    \item \textbf{Sezione 1 – Introduzione}: scopo, definizioni e panoramica.
    \item \textbf{Sezione 2 – Descrizione Generale}: contesto, assunzioni e dipendenze.
    \item \textbf{Sezione 3 – Requisiti Funzionali}: funzionalità dettagliate del sistema.
    \item \textbf{Sezione 4 – Requisiti Non Funzionali}: qualità, performance, sicurezza.
    \item \textbf{Sezione 5 – Requisiti di Sistema e Vincoli Tecnologici}.
    \item \textbf{Sezione 6 – Modello dei Casi d'Uso}: attori e use case principali.
    \item \textbf{Sezione 7 – Schema del Database}.
    \item \textbf{Sezione 8 – Requisiti di Consegna e Valutazione}.
    \item \textbf{Sezione 9 – Appendici}.
\end{itemize}

% ============================================================
\newpage
\section{Descrizione Generale}

\subsection{Prospettiva del Sistema}

\textit{GuessTheWord} è un sistema software \textbf{autonomo e auto-contenuto}, progettato per funzionare interamente in locale (localhost) o su rete LAN. Non dipende da servizi cloud, API esterne o infrastrutture di terze parti.

Il sistema segue un'architettura \textbf{client–server} classica:

\begin{center}
\begin{tcolorbox}[colback=lightBlue, colframe=primaryBlue, width=0.9\linewidth]
\centering
\texttt{[Client JavaFX 1]} $\longleftrightarrow$ \texttt{[Server JavaFX]} $\longleftrightarrow$ \texttt{[Client JavaFX 2]}\\
\vspace{0.3cm}
\texttt{[Server JavaFX]} $\longleftrightarrow$ \texttt{[Database SQLite]}\\
\texttt{[Server JavaFX]} $\longleftrightarrow$ \texttt{[File .properties, File TXT, Dati Serializzati]}
\end{tcolorbox}
\end{center}

La comunicazione tra client e server avviene tramite \textbf{socket TCP}. Il server è l'unico componente che accede al database e che conosce le risposte corrette alle sfide.

\subsection{Funzioni Principali del Sistema}

Il sistema svolge le seguenti macro-funzioni:

\begin{enumerate}[leftmargin=*]
    \item \textbf{Gestione utenti}: registrazione, autenticazione, distinzione di ruolo (admin/giocatore).
    \item \textbf{Analisi documenti}: caricamento di file TXT, calcolo della Term Frequency con Stream API, salvataggio serializzato dei risultati.
    \item \textbf{Preparazione sfide}: selezione di estratti testuali, cifratura con Cesare, invio ai client.
    \item \textbf{Gestione partite}: coordinamento in tempo reale, validazione risposte, dichiarazione vincitore.
    \item \textbf{Persistenza e classifiche}: storico sfide, statistiche per utente, consultazione classifiche.
    \item \textbf{Configurazione}: lettura da file \code{.properties} per IP, porta e altri parametri.
\end{enumerate}

\subsection{Caratteristiche degli Utenti}

\begin{longtable}{|p{3cm}|p{5cm}|p{5.5cm}|}
    \hline
    \rowcolor{primaryBlue!20}
    \textbf{Ruolo} & \textbf{Descrizione} & \textbf{Competenze Richieste} \\
    \hline
    \endhead
    \textbf{Amministratore} &
    Utente speciale con accesso all'interfaccia di gestione del server. Gestisce i documenti e monitora il sistema. &
    Conoscenza di base del sistema; le sue credenziali sono pre-caricate nel DB. \\
    \hline
    \textbf{Giocatore} &
    Utente che si registra o accede tramite il client e partecipa alle sfide. &
    Nessuna competenza tecnica richiesta; interazione esclusivamente tramite GUI. \\
    \hline
\end{longtable}

\subsection{Vincoli Generali}

\begin{constraintbox}
\begin{itemize}[leftmargin=*]
    \item Il sistema \mustshall{} essere compatibile con \textbf{JDK 8}.
    \item Il sistema \mustshall{} utilizzare esclusivamente le tecnologie obbligatorie elencate nella Sezione 5.
    \item Tutti i percorsi di file e database \mustshall{} essere \textbf{relativi}, mai assoluti.
    \item Nessuna dipendenza esterna \mustshall{} essere richiesta al di fuori di quelle incluse nell'archivio di consegna.
    \item Il sistema \mustshall{} essere avviabile tramite semplice esecuzione di file JAR.
\end{itemize}
\end{constraintbox}

\subsection{Assunzioni e Dipendenze}

\begin{itemize}[leftmargin=*]
    \item Si assume che il sistema operativo host supporti Java 8 e JavaFX 8.
    \item Si assume che il server e i client vengano avviati sulla stessa macchina (localhost) o in rete LAN durante i test.
    \item Si assume che il database SQLite venga inizializzato automaticamente all'avvio con lo schema e gli account predefiniti.
    \item Si assume che i file TXT da analizzare siano già disponibili nel filesystem prima dell'avvio del server.
    \item La disponibilità della rete è necessaria per la comunicazione via socket tra client e server.
\end{itemize}

% ============================================================
\newpage
\section{Requisiti Funzionali}

I requisiti sono classificati con il seguente schema di priorità:
\begin{itemize}[leftmargin=*]
    \item \textbf{[MUST]}: requisito obbligatorio, indispensabile per il funzionamento del sistema.
    \item \textbf{[SHOULD]}: requisito fortemente consigliato.
    \item \textbf{[MAY]}: requisito opzionale o facoltativo.
\end{itemize}

% ── 3.1 Server ─────────────────────────────────────────────
\subsection{Requisiti del Server}

\subsubsection{Configurazione e Avvio}

\begin{reqbox}{RF-SRV-01 \quad Lettura configurazione da file \quad [MUST]}
Il server \mustshall{} leggere la porta di ascolto da un file di configurazione denominato \code{server.properties}, nella forma:
\begin{lstlisting}[language={}]
server.port=5000
\end{lstlisting}
Il file \mustshall{} essere localizzabile tramite percorso relativo alla directory di esecuzione.
\end{reqbox}

\begin{reqbox}{RF-SRV-02 \quad Apertura ServerSocket \quad [MUST]}
Il server \mustshall{} aprire un \code{ServerSocket} sulla porta specificata nel file di configurazione e rimanere in ascolto di connessioni client in modo continuo.
\end{reqbox}

\begin{reqbox}{RF-SRV-03 \quad Interfaccia grafica amministratore \quad [MUST]}
Il server \mustshall{} fornire un'interfaccia grafica JavaFX (con file FXML) accessibile all'amministratore, che consenta di:
\begin{itemize}
    \item selezionare uno o più file di testo (\code{.txt}) da analizzare;
    \item avviare l'analisi del contenuto;
    \item salvare i risultati dell'analisi in formato serializzato;
    \item caricare risultati di analisi precedentemente salvati.
\end{itemize}
\end{reqbox}

\subsubsection{Autenticazione e Gestione Utenti}

\begin{reqbox}{RF-SRV-04 \quad Autenticazione client \quad [MUST]}
Il server \mustshall{} ricevere le credenziali (username e password) inviate dai client, verificarle nel database SQLite e rispondere con l'esito dell'autenticazione.
\end{reqbox}

\begin{reqbox}{RF-SRV-05 \quad Registrazione nuovi utenti \quad [MUST]}
Il server \mustshall{} ricevere richieste di registrazione dai client, verificare l'unicità dello username e, in caso di successo, inserire il nuovo utente nel database con il ruolo di \textit{giocatore}.
\end{reqbox}

\begin{reqbox}{RF-SRV-06 \quad Gestione account amministratore \quad [MUST]}
Il server \mustshall{} riconoscere l'account amministratore e consentirne l'accesso all'interfaccia di gestione. Le credenziali dell'amministratore \mustshall{} essere pre-caricate nel database.
\end{reqbox}

\subsubsection{Analisi Documenti}

\begin{reqbox}{RF-SRV-07 \quad Analisi Term Frequency \quad [MUST]}
Il server \mustshall{} elaborare i documenti TXT selezionati calcolando la \textbf{Term Frequency} (frequenza relativa) di ciascuna parola. L'elaborazione \mustshall{} utilizzare la \textbf{Stream API di Java}.
\end{reqbox}

\begin{reqbox}{RF-SRV-08 \quad Esecuzione asincrona dell'analisi \quad [MUST]}
L'analisi dei documenti \mustshall{} essere eseguita in modo asincrono tramite \textbf{JavaFX Service / Task}, in modo che l'interfaccia grafica del server rimanga reattiva durante l'elaborazione.
\end{reqbox}

\begin{reqbox}{RF-SRV-09 \quad Salvataggio serializzato dei risultati \quad [MUST]}
Il server \mustshall{} consentire il salvataggio dei risultati dell'analisi (mappa parola $\rightarrow$ frequenza) in formato \textbf{serializzato} su file, e il successivo ricaricamento all'avvio, evitando di rieseguire l'analisi.
\end{reqbox}

\subsubsection{Preparazione e Gestione della Sfida}

\begin{reqbox}{RF-SRV-10 \quad Attesa di due client connessi e autenticati \quad [MUST]}
Il server \mustshall{} avviare una sfida solamente quando \textbf{entrambi i client} si sono connessi e hanno completato con successo l'autenticazione.
\end{reqbox}

\begin{reqbox}{RF-SRV-11 \quad Selezione dell'estratto testuale \quad [MUST]}
Il server \mustshall{} selezionare un estratto di testo da uno dei documenti analizzati. La selezione dell'estratto \should{} avvenire in modo da garantire varietà tra le sfide.
\end{reqbox}

\begin{reqbox}{RF-SRV-12 \quad Applicazione del cifrario di Cesare \quad [MUST]}
Il server \mustshall{} selezionare una o più parole dall'estratto e sostituirle con la rispettiva versione cifrata tramite il \textbf{cifrario di Cesare}, applicando uno spostamento scelto \textbf{casualmente}. Le versioni cifrate \mustshall{} essere chiaramente distinguibili nel testo inviato ai client.
\end{reqbox}

\begin{reqbox}{RF-SRV-13 \quad Invio del testo cifrato ai client \quad [MUST]}
Il server \mustshall{} inviare a entrambi i client:
\begin{itemize}
    \item il testo dell'estratto con le parole cifrate;
    \item la durata del conto alla rovescia (timeout della sfida).
\end{itemize}
\end{reqbox}

\begin{reqbox}{RF-SRV-14 \quad Ricezione e validazione delle risposte \quad [MUST]}
Il server \mustshall{} ricevere le risposte inviate dai client, validarle confrontandole con la parola originale (in modo case-insensitive o secondo le regole stabilite) e determinare il vincitore come il \textbf{primo giocatore che risponde correttamente}.
\end{reqbox}

\begin{reqbox}{RF-SRV-15 \quad Notifica dell'esito \quad [MUST]}
Il server \mustshall{} notificare \textbf{entrambi} i client dell'esito della sfida, indicando chi ha vinto e, se applicabile, la risposta corretta.
\end{reqbox}

\begin{reqbox}{RF-SRV-16 \quad Registrazione dell'esito nel database \quad [MUST]}
Al termine di ogni sfida, il server \mustshall{} registrare nel database:
\begin{itemize}
    \item gli utenti partecipanti;
    \item l'esito (vincitore e perdente);
    \item il riferimento temporale (data e ora della sfida).
\end{itemize}
\end{reqbox}

\begin{reqbox}{RF-SRV-17 \quad Livelli di difficoltà (opzionale) \quad [MAY]}
Il server \may{} implementare livelli di difficoltà variabili per le sfide, basati su:
\begin{itemize}
    \item frequenza della parola selezionata (parola più rara = difficoltà maggiore);
    \item lunghezza della parola;
    \item numero di parole cifrate simultaneamente;
    \item valore dello shift del cifrario.
\end{itemize}
\end{reqbox}

\subsubsection{Classifiche e Statistiche (Lato Server)}

\begin{reqbox}{RF-SRV-18 \quad Visualizzazione classifiche dall'interfaccia admin \quad [SHOULD]}
Il server \should{} consentire all'amministratore di visualizzare classifiche e statistiche sulle partite giocate tramite l'interfaccia grafica, inclusi:
\begin{itemize}
    \item numero di vittorie per utente;
    \item numero totale di partite disputate per utente;
    \item tempo medio di risposta.
\end{itemize}
\end{reqbox}

% ── 3.2 Client ─────────────────────────────────────────────
\subsection{Requisiti del Client}

\subsubsection{Configurazione}

\begin{reqbox}{RF-CLT-01 \quad Lettura configurazione da file \quad [MUST]}
Ciascun client \mustshall{} leggere l'indirizzo IP e la porta del server da un file di configurazione denominato \code{client.properties}, nella forma:
\begin{lstlisting}[language={}]
server.ip=127.0.0.1
server.port=5000
\end{lstlisting}
\end{reqbox}

\subsubsection{Autenticazione}

\begin{reqbox}{RF-CLT-02 \quad Interfaccia di login e registrazione \quad [MUST]}
Il client \mustshall{} presentare una finestra iniziale (JavaFX/FXML) con:
\begin{itemize}
    \item campo di inserimento per lo username;
    \item campo di inserimento per la password;
    \item pulsante \textit{Login};
    \item pulsante \textit{Registrazione}.
\end{itemize}
\end{reqbox}

\begin{reqbox}{RF-CLT-03 \quad Comunicazione credenziali al server \quad [MUST]}
Il client \mustshall{} trasmettere al server le credenziali inserite dall'utente e riceverne l'esito, informando l'utente in caso di errore (credenziali errate, username già esistente, ecc.).
\end{reqbox}

\begin{reqbox}{RF-CLT-04 \quad Attesa dell'avversario \quad [MUST]}
Dopo l'autenticazione, il client \mustshall{} mostrare uno stato di \textbf{attesa} finché l'altro sfidante non si connette e autentica.
\end{reqbox}

\subsubsection{Finestra di Gioco}

\begin{reqbox}{RF-CLT-05 \quad Visualizzazione del testo cifrato \quad [MUST]}
Durante la sfida, il client \mustshall{} visualizzare il testo ricevuto dal server con le parole cifrate \textbf{chiaramente evidenziate} o distinguibili dal resto del testo.
\end{reqbox}

\begin{reqbox}{RF-CLT-06 \quad Inserimento e invio della risposta \quad [MUST]}
Il client \mustshall{} fornire:
\begin{itemize}
    \item un campo di testo per l'inserimento della risposta;
    \item un pulsante per l'invio della risposta al server.
\end{itemize}
\end{reqbox}

\begin{reqbox}{RF-CLT-07 \quad Timer visibile \quad [MUST]}
Il client \mustshall{} mostrare un \textbf{conto alla rovescia} visibile, sincronizzato con la durata della sfida comunicata dal server.
\end{reqbox}

\begin{reqbox}{RF-CLT-08 \quad Indicatore di stato della partita \quad [MUST]}
Il client \mustshall{} mostrare lo stato attuale della partita, che include almeno i seguenti stati:
\begin{itemize}
    \item \textit{In attesa} — connesso, in attesa dell'avversario;
    \item \textit{In gioco} — sfida in corso;
    \item \textit{Terminata} — sfida conclusa, esito ricevuto.
\end{itemize}
\end{reqbox}

\begin{reqbox}{RF-CLT-09 \quad Visualizzazione esito \quad [MUST]}
Il client \mustshall{} visualizzare l'esito della sfida ricevuto dal server, indicando se l'utente ha vinto o perso e la risposta corretta.
\end{reqbox}

\subsubsection{Storico e Classifiche (Lato Client)}

\begin{reqbox}{RF-CLT-10 \quad Visualizzazione storico personale \quad [SHOULD]}
Il client \should{} consentire a ciascun utente di visualizzare lo storico delle proprie sfide precedenti, con esiti e date.
\end{reqbox}

% ── 3.3 Persistenza ────────────────────────────────────────
\subsection{Requisiti di Persistenza dei Dati}

\begin{reqbox}{RF-DB-01 \quad Tabella Utenti \quad [MUST]}
Il database \mustshall{} contenere una struttura dati per gli utenti, che includa almeno:
\begin{itemize}
    \item identificativo univoco dell'utente;
    \item username (univoco);
    \item password (si raccomanda hashing);
    \item ruolo (amministratore / giocatore).
\end{itemize}
\end{reqbox}

\begin{reqbox}{RF-DB-02 \quad Tabella Sfide \quad [MUST]}
Il database \mustshall{} contenere una struttura dati per le sfide disputate, che includa almeno:
\begin{itemize}
    \item identificativo univoco della sfida;
    \item riferimento temporale (data e ora);
    \item testo dell'estratto e parola/e cifrata/e (opzionale).
\end{itemize}
\end{reqbox}

\begin{reqbox}{RF-DB-03 \quad Tabella Risultati \quad [MUST]}
Il database \mustshall{} contenere una struttura dati per i risultati, che includa almeno:
\begin{itemize}
    \item riferimento alla sfida;
    \item riferimento agli utenti partecipanti;
    \item esito per ciascun utente (vittoria / sconfitta);
    \item tempo di risposta (opzionale).
\end{itemize}
\end{reqbox}

\begin{reqbox}{RF-DB-04 \quad Statistiche estraibili \quad [MUST]}
Il sistema \mustshall{} supportare l'estrazione delle seguenti statistiche tramite query sul database:
\begin{itemize}
    \item numero di vittorie per utente;
    \item numero di partite disputate per utente;
    \item tempo medio di risposta per utente.
\end{itemize}
\end{reqbox}

\begin{reqbox}{RF-DB-05 \quad Account predefiniti \quad [MUST]}
Il database \mustshall{} contenere, all'avvio del sistema, almeno i seguenti account pre-caricati:
\begin{itemize}
    \item \textbf{1 account amministratore};
    \item \textbf{2 account giocatore}.
\end{itemize}
Le credenziali di tali account \mustshall{} essere documentate nella relazione e nel file \code{readme.txt}.
\end{reqbox}

% ============================================================
\newpage
\section{Requisiti Non Funzionali}

\subsection{Usabilità}

\begin{reqbox}{RNF-USA-01 \quad Interfaccia intuitiva \quad [MUST]}
Le interfacce grafiche (server e client) \mustshall{} essere intuitive e navigabili senza necessità di documentazione aggiuntiva per un utente con minime competenze informatiche.
\end{reqbox}

\begin{reqbox}{RNF-USA-02 \quad Feedback visivo \quad [MUST]}
Il sistema \mustshall{} fornire feedback visivo per ogni azione significativa dell'utente (login, invio risposta, fine partita, ecc.).
\end{reqbox}

\begin{reqbox}{RNF-USA-03 \quad Styling CSS (opzionale) \quad [MAY]}
Le interfacce JavaFX \may{} essere stilizzate tramite fogli di stile CSS per migliorare l'aspetto grafico.
\end{reqbox}

\subsection{Affidabilità}

\begin{reqbox}{RNF-AFF-01 \quad Gestione disconnessioni \quad [SHOULD]}
Il server \should{} gestire in modo robusto la disconnessione inattesa di un client durante una sfida, notificando l'altro client e aggiornando lo stato della partita nel database.
\end{reqbox}

\begin{reqbox}{RNF-AFF-02 \quad Validazione input \quad [MUST]}
Il sistema \mustshall{} validare tutti gli input utente (campi vuoti, caratteri non validi) sia lato client (prima dell'invio) sia lato server (prima dell'elaborazione).
\end{reqbox}

\subsection{Performance}

\begin{reqbox}{RNF-PRF-01 \quad Reattività dell'interfaccia \quad [MUST]}
L'interfaccia grafica del server \mustshall{} rimanere reattiva durante l'analisi dei documenti, grazie all'uso di \textbf{JavaFX Service / Task}.
\end{reqbox}

\begin{reqbox}{RNF-PRF-02 \quad Latenza di comunicazione \quad [SHOULD]}
La latenza di comunicazione tra client e server, in condizioni di rete locale (localhost / LAN), \should{} essere sufficientemente bassa da non compromettere l'esperienza di gioco.
\end{reqbox}

\subsection{Manutenibilità}

\begin{reqbox}{RNF-MNT-01 \quad Documentazione Javadoc \quad [MUST]}
Il codice sorgente \mustshall{} essere documentato con \textbf{Javadoc} per tutte le classi e i metodi pubblici.
\end{reqbox}

\begin{reqbox}{RNF-MNT-02 \quad Configurabilità \quad [MUST]}
Tutti i parametri configurabili (porta, IP, ecc.) \mustshall{} essere gestiti tramite file \code{.properties} e non codificati direttamente nel sorgente (\textit{no hard-coding}).
\end{reqbox}

\subsection{Portabilità}

\begin{reqbox}{RNF-PRT-01 \quad Compatibilità JDK 8 \quad [MUST]}
Tutti i file JAR prodotti \mustshall{} essere compatibili con \textbf{JDK 8} e non richiedere versioni successive.
\end{reqbox}

\begin{reqbox}{RNF-PRT-02 \quad Dipendenze auto-contenute \quad [MUST]}
L'archivio di consegna \mustshall{} includere tutte le dipendenze necessarie. Non \mustshall{} essere richiesta nessuna installazione aggiuntiva da parte del revisore.
\end{reqbox}

\begin{reqbox}{RNF-PRT-03 \quad Percorsi relativi \quad [MUST]}
Tutti i percorsi a file, database e risorse \mustshall{} essere \textbf{relativi} alla directory di esecuzione, mai assoluti.
\end{reqbox}

% ============================================================
\newpage
\section{Requisiti di Sistema e Vincoli Tecnologici}

\subsection{Tecnologie Obbligatorie}

Le seguenti tecnologie \mustshall{} essere utilizzate nel progetto:

\begin{longtable}{|p{4cm}|p{9cm}|}
    \hline
    \rowcolor{primaryBlue!20}
    \textbf{Tecnologia} & \textbf{Utilizzo nel Progetto} \\
    \hline
    \endhead
    \textbf{Java (JDK 8)} & Linguaggio principale per server e client \\
    \hline
    \textbf{JavaFX + FXML} & Interfaccia grafica per server (admin) e client (giocatori) \\
    \hline
    \textbf{Socket TCP} & Comunicazione client–server in tempo reale \\
    \hline
    \textbf{Stream API} & Elaborazione del testo per l'analisi della Term Frequency \\
    \hline
    \textbf{JavaFX Service / Task} & Esecuzione asincrona dell'analisi documenti \\
    \hline
    \textbf{SQLite (via JDBC)} & Persistenza di utenti, sfide e risultati \\
    \hline
    \textbf{File .properties} & Configurazione di server e client \\
    \hline
    \textbf{Serializzazione Java} & Salvataggio e ricaricamento dei risultati dell'analisi \\
    \hline
\end{longtable}

\subsection{Tecnologie Facoltative}

\begin{longtable}{|p{4cm}|p{9cm}|}
    \hline
    \rowcolor{primaryBlue!20}
    \textbf{Tecnologia} & \textbf{Utilizzo nel Progetto} \\
    \hline
    \endhead
    \textbf{CSS (JavaFX)} & Styling delle interfacce grafiche \\
    \hline
\end{longtable}

\subsection{Ambiente di Esecuzione}

\begin{itemize}[leftmargin=*]
    \item Sistema Operativo: qualsiasi sistema supportato da JDK 8 (Windows, macOS, Linux).
    \item JDK 8 o superiore installato sulla macchina del revisore.
    \item Accesso alla rete locale (socket su localhost o LAN).
\end{itemize}

\subsection{Struttura dell'Archivio di Esecuzione}

L'archivio \code{eseguibili.zip} \mustshall{} avere la seguente struttura minima:

\begin{lstlisting}[language={}]
eseguibili/
├── server.jar
├── client.jar
├── readme.txt
├── properties/
│   ├── server.properties
│   └── client.properties
├── db/
│   └── database.db
├── documents/
└── data/
\end{lstlisting}

I file possono essere organizzati in modo diverso purché tutti i componenti necessari siano presenti e il sistema sia immediatamente eseguibile.

% ============================================================
\newpage
\section{Modello dei Casi d'Uso}

\subsection{Attori del Sistema}

\begin{longtable}{|p{3cm}|p{10cm}|}
    \hline
    \rowcolor{primaryBlue!20}
    \textbf{Attore} & \textbf{Descrizione} \\
    \hline
    \endhead
    \textbf{Amministratore} & Utente speciale che gestisce il server tramite l'interfaccia admin. Carica i documenti, avvia le analisi e monitora le classifiche. \\
    \hline
    \textbf{Giocatore} & Utente registrato che si connette tramite il client per partecipare alle sfide. \\
    \hline
    \textbf{Server} & Componente software attivo che coordina le connessioni, le sfide e la persistenza. Attore secondario nei casi d'uso del client. \\
    \hline
\end{longtable}

\subsection{Casi d'Uso — Amministratore}

\begin{longtable}{|p{1.5cm}|p{3.5cm}|p{8cm}|}
    \hline
    \rowcolor{primaryBlue!20}
    \textbf{ID} & \textbf{Nome} & \textbf{Descrizione} \\
    \hline
    \endhead
    UC-A01 & Login amministratore & L'amministratore accede al server con le proprie credenziali tramite l'interfaccia admin. \\
    \hline
    UC-A02 & Caricamento documenti & L'amministratore seleziona uno o più file TXT da filesystem per l'analisi. \\
    \hline
    UC-A03 & Avvio analisi & L'amministratore avvia l'analisi TF dei documenti caricati; la UI rimane reattiva durante l'elaborazione. \\
    \hline
    UC-A04 & Salvataggio analisi & L'amministratore salva i risultati dell'analisi in formato serializzato. \\
    \hline
    UC-A05 & Caricamento analisi & L'amministratore carica i risultati di un'analisi precedentemente salvata. \\
    \hline
    UC-A06 & Visualizzazione classifiche & L'amministratore consulta le statistiche e le classifiche degli utenti. \\
    \hline
\end{longtable}

\subsection{Casi d'Uso — Giocatore}

\begin{longtable}{|p{1.5cm}|p{3.5cm}|p{8cm}|}
    \hline
    \rowcolor{primaryBlue!20}
    \textbf{ID} & \textbf{Nome} & \textbf{Descrizione} \\
    \hline
    \endhead
    UC-G01 & Registrazione & Il giocatore crea un nuovo account inserendo username e password. Il server verifica l'unicità e registra l'utente. \\
    \hline
    UC-G02 & Login & Il giocatore accede con le proprie credenziali. Il server autentica e conferma l'accesso. \\
    \hline
    UC-G03 & Attesa partita & Dopo il login, il giocatore attende che il secondo client si connetta e autentichi. \\
    \hline
    UC-G04 & Partecipazione sfida & Il giocatore riceve il testo cifrato, lo analizza e invia la propria risposta entro il tempo limite. \\
    \hline
    UC-G05 & Visualizzazione esito & Al termine della sfida, il giocatore visualizza se ha vinto o perso e la parola corretta. \\
    \hline
    UC-G06 & Visualizzazione storico & Il giocatore consulta lo storico delle proprie sfide precedenti con relativi esiti. \\
    \hline
\end{longtable}

\subsection{Caso d'Uso Dettagliato — UC-G04: Partecipazione alla Sfida}

\begin{tcolorbox}[colback=lightGray, colframe=darkGray!60, title=\textbf{UC-G04 — Partecipazione alla Sfida}]
\textbf{Attore Primario:} Giocatore\\
\textbf{Attori Secondari:} Server, secondo Giocatore\\[0.3cm]
\textbf{Pre-condizioni:}
\begin{itemize}
    \item Entrambi i giocatori sono autenticati e connessi al server.
    \item Il server ha almeno un documento analizzato disponibile.
\end{itemize}
\textbf{Flusso Principale:}
\begin{enumerate}
    \item Il server seleziona un estratto da un documento analizzato.
    \item Il server sceglie una o più parole e le cifra con Cesare (shift casuale).
    \item Il server invia il testo cifrato e il timeout a entrambi i client.
    \item I client visualizzano il testo con il timer attivo.
    \item I giocatori inseriscono e inviano la loro risposta.
    \item Il server riceve le risposte e valida.
    \item Il primo giocatore con risposta corretta viene dichiarato vincitore.
    \item Il server notifica entrambi i client e registra l'esito nel database.
\end{enumerate}
\textbf{Flusso Alternativo A — Timeout:}
\begin{itemize}
    \item Se nessun giocatore risponde correttamente entro il timeout, la sfida termina in pareggio o senza vincitore.
    \item Il server notifica entrambi i client e registra l'esito.
\end{itemize}
\textbf{Post-condizioni:}
\begin{itemize}
    \item L'esito della sfida è registrato nel database.
    \item Entrambi i client mostrano il risultato finale.
\end{itemize}
\end{tcolorbox}

% ============================================================
\newpage
\section{Schema del Database}

Il design del database è lasciato alla libera progettazione del gruppo. La struttura seguente rappresenta un modello di riferimento che soddisfa i requisiti descritti nel documento.

\subsection{Tabella: \texttt{users}}

\begin{longtable}{|p{3.5cm}|p{2.5cm}|p{7cm}|}
    \hline
    \rowcolor{primaryBlue!20}
    \textbf{Colonna} & \textbf{Tipo} & \textbf{Descrizione} \\
    \hline
    \endhead
    \code{user\_id} & INTEGER PK & Chiave primaria auto-incrementale \\
    \hline
    \code{username} & TEXT UNIQUE & Nome utente univoco \\
    \hline
    \code{password\_hash} & TEXT & Hash della password \\
    \hline
    \code{role} & TEXT & Ruolo: \code{admin} o \code{player} \\
    \hline
    \code{created\_at} & DATETIME & Data e ora di registrazione \\
    \hline
\end{longtable}

\subsection{Tabella: \texttt{challenges}}

\begin{longtable}{|p{3.5cm}|p{2.5cm}|p{7cm}|}
    \hline
    \rowcolor{primaryBlue!20}
    \textbf{Colonna} & \textbf{Tipo} & \textbf{Descrizione} \\
    \hline
    \endhead
    \code{challenge\_id} & INTEGER PK & Chiave primaria auto-incrementale \\
    \hline
    \code{played\_at} & DATETIME & Data e ora della sfida \\
    \hline
    \code{text\_excerpt} & TEXT & Estratto di testo usato \\
    \hline
    \code{encrypted\_word} & TEXT & Parola/e cifrata/e \\
    \hline
    \code{original\_word} & TEXT & Parola/e originale/i \\
    \hline
    \code{caesar\_shift} & INTEGER & Spostamento del cifrario applicato \\
    \hline
\end{longtable}

\subsection{Tabella: \texttt{results}}

\begin{longtable}{|p{3.5cm}|p{2.5cm}|p{7cm}|}
    \hline
    \rowcolor{primaryBlue!20}
    \textbf{Colonna} & \textbf{Tipo} & \textbf{Descrizione} \\
    \hline
    \endhead
    \code{result\_id} & INTEGER PK & Chiave primaria auto-incrementale \\
    \hline
    \code{challenge\_id} & INTEGER FK & Riferimento alla sfida \\
    \hline
    \code{user\_id} & INTEGER FK & Riferimento all'utente \\
    \hline
    \code{outcome} & TEXT & Esito: \code{win} / \code{loss} / \code{draw} \\
    \hline
    \code{response\_time\_ms} & INTEGER & Tempo di risposta in millisecondi \\
    \hline
\end{longtable}

\subsection{Query di Esempio per le Statistiche}

\begin{lstlisting}[language=SQL, caption={Numero di vittorie per utente}]
SELECT u.username, COUNT(*) AS wins
FROM results r
JOIN users u ON r.user_id = u.user_id
WHERE r.outcome = 'win'
GROUP BY r.user_id
ORDER BY wins DESC;
\end{lstlisting}

\begin{lstlisting}[language=SQL, caption={Tempo medio di risposta per utente}]
SELECT u.username,
       COUNT(*) AS games_played,
       AVG(r.response_time_ms) AS avg_response_ms
FROM results r
JOIN users u ON r.user_id = u.user_id
GROUP BY r.user_id
ORDER BY avg_response_ms ASC;
\end{lstlisting}

% ============================================================
\newpage
\section{Requisiti di Consegna e Valutazione}

\subsection{Materiale Obbligatorio}

\begin{reqbox}{RC-01 \quad Codice sorgente \quad [MUST]}
\mustshall{} essere consegnato il codice sorgente completo del progetto (moduli server e client), corredato di \textbf{Javadoc} per tutte le classi e i metodi pubblici.
\end{reqbox}

\begin{reqbox}{RC-02 \quad Relazione tecnica \quad [MUST]}
\mustshall{} essere prodotta una relazione sintetica che documenti:
\begin{itemize}
    \item Architettura del sistema;
    \item Principali scelte progettuali;
    \item Specifica dei requisiti (Use Cases, Mockup interfacce);
    \item Progettazione (Class Diagram, Sequence Diagram, Modello E-R);
    \item Tecnologie utilizzate.
\end{itemize}
\end{reqbox}

\begin{reqbox}{RC-03 \quad Presentazione \quad [MUST]}
\mustshall{} essere prodotta una \textbf{presentazione PowerPoint} sintetica del progetto.
\end{reqbox}

\begin{reqbox}{RC-04 \quad Archivio eseguibili \quad [MUST]}
\mustshall{} essere consegnato un archivio ZIP denominato \code{eseguibili.zip} contenente:
\begin{itemize}
    \item \code{server.jar}
    \item \code{client.jar}
    \item \code{readme.txt} (con credenziali degli account predefiniti)
    \item File di configurazione (\code{server.properties}, \code{client.properties})
    \item Database SQLite inizializzato (\code{database.db})
    \item Tutte le dipendenze necessarie
\end{itemize}
\end{reqbox}

\subsection{Scadenza}

\begin{constraintbox}
La consegna del progetto \mustshall{} avvenire \textbf{7 giorni prima dell'appello d'esame}.
\end{constraintbox}

\subsection{Requisiti per l'Archivio Eseguibile}

\begin{reqbox}{RC-05 \quad Avvio immediato \quad [MUST]}
Il revisore \mustshall{} poter avviare il sistema con i seguenti comandi, senza ulteriori configurazioni:
\begin{lstlisting}[language={}]
java -jar server.jar
java -jar client.jar
java -jar client.jar
\end{lstlisting}
\end{reqbox}

\begin{reqbox}{RC-06 \quad Account predefiniti per il testing \quad [MUST]}
Il database consegnato \mustshall{} contenere:
\begin{itemize}
    \item \textbf{1 account amministratore} (per accesso all'interfaccia admin del server);
    \item \textbf{2 account giocatore} (per testare immediatamente una sfida completa).
\end{itemize}
Le credenziali \mustshall{} essere riportate sia nella relazione sia nel file \code{readme.txt}.
\end{reqbox}

\subsection{Modalità di Valutazione — Esame Finale}

\begin{itemize}[leftmargin=*]
    \item L'esame consiste in una \textbf{presentazione di gruppo} (circa 10 minuti), comprensiva di \textbf{demo live} del progetto.
    \item Segue un \textbf{colloquio individuale} sulle scelte progettuali e sugli argomenti del corso.
    \item Tutti i componenti del gruppo \mustshall{} dimostrare padronanza dell'intero progetto.
    \item Il colloquio individuale non è limitato al progetto: può riguardare qualsiasi argomento trattato durante il corso.
    \item La consegna di archivi incompleti o non eseguibili influisce negativamente sulla valutazione.
    \item È preferibile che tutti i membri del gruppo sostengano l'esame nello stesso appello.
    \item Gli studenti che rimandano l'esame \mustshall{} ripresentare il progetto nell'appello successivo.
\end{itemize}

% ============================================================
\newpage
\section{Appendici}

\subsection{Appendice A — Riepilogo dei Requisiti}

\subsubsection{Requisiti Funzionali}

\begin{longtable}{|p{2.5cm}|p{7cm}|p{2cm}|p{1.5cm}|}
    \hline
    \rowcolor{primaryBlue!20}
    \textbf{ID} & \textbf{Descrizione sintetica} & \textbf{Componente} & \textbf{Priorità} \\
    \hline
    \endhead
    RF-SRV-01 & Lettura porta da server.properties & Server & MUST \\
    \hline
    RF-SRV-02 & Apertura ServerSocket & Server & MUST \\
    \hline
    RF-SRV-03 & Interfaccia admin JavaFX & Server & MUST \\
    \hline
    RF-SRV-04 & Autenticazione client & Server & MUST \\
    \hline
    RF-SRV-05 & Registrazione nuovi utenti & Server & MUST \\
    \hline
    RF-SRV-06 & Gestione account admin & Server & MUST \\
    \hline
    RF-SRV-07 & Analisi TF con Stream API & Server & MUST \\
    \hline
    RF-SRV-08 & Analisi asincrona con Service/Task & Server & MUST \\
    \hline
    RF-SRV-09 & Serializzazione risultati analisi & Server & MUST \\
    \hline
    RF-SRV-10 & Attesa due client autenticati & Server & MUST \\
    \hline
    RF-SRV-11 & Selezione estratto testuale & Server & MUST \\
    \hline
    RF-SRV-12 & Cifrario di Cesare con shift casuale & Server & MUST \\
    \hline
    RF-SRV-13 & Invio testo cifrato ai client & Server & MUST \\
    \hline
    RF-SRV-14 & Validazione risposte & Server & MUST \\
    \hline
    RF-SRV-15 & Notifica esito a entrambi i client & Server & MUST \\
    \hline
    RF-SRV-16 & Registrazione esito nel DB & Server & MUST \\
    \hline
    RF-SRV-17 & Livelli di difficoltà & Server & MAY \\
    \hline
    RF-SRV-18 & Classifiche nell'interfaccia admin & Server & SHOULD \\
    \hline
    RF-CLT-01 & Lettura IP/porta da client.properties & Client & MUST \\
    \hline
    RF-CLT-02 & Interfaccia login/registrazione & Client & MUST \\
    \hline
    RF-CLT-03 & Comunicazione credenziali al server & Client & MUST \\
    \hline
    RF-CLT-04 & Attesa avversario & Client & MUST \\
    \hline
    RF-CLT-05 & Visualizzazione testo cifrato & Client & MUST \\
    \hline
    RF-CLT-06 & Inserimento e invio risposta & Client & MUST \\
    \hline
    RF-CLT-07 & Timer visibile & Client & MUST \\
    \hline
    RF-CLT-08 & Indicatore stato partita & Client & MUST \\
    \hline
    RF-CLT-09 & Visualizzazione esito & Client & MUST \\
    \hline
    RF-CLT-10 & Storico personale & Client & SHOULD \\
    \hline
    RF-DB-01 & Struttura tabella Utenti & Database & MUST \\
    \hline
    RF-DB-02 & Struttura tabella Sfide & Database & MUST \\
    \hline
    RF-DB-03 & Struttura tabella Risultati & Database & MUST \\
    \hline
    RF-DB-04 & Statistiche estraibili & Database & MUST \\
    \hline
    RF-DB-05 & Account predefiniti nel DB & Database & MUST \\
    \hline
\end{longtable}

\subsubsection{Requisiti Non Funzionali}

\begin{longtable}{|p{2.5cm}|p{7cm}|p{2.5cm}|p{1.5cm}|}
    \hline
    \rowcolor{primaryBlue!20}
    \textbf{ID} & \textbf{Descrizione sintetica} & \textbf{Categoria} & \textbf{Priorità} \\
    \hline
    \endhead
    RNF-USA-01 & Interfaccia intuitiva & Usabilità & MUST \\
    \hline
    RNF-USA-02 & Feedback visivo & Usabilità & MUST \\
    \hline
    RNF-USA-03 & Styling CSS & Usabilità & MAY \\
    \hline
    RNF-AFF-01 & Gestione disconnessioni & Affidabilità & SHOULD \\
    \hline
    RNF-AFF-02 & Validazione input & Affidabilità & MUST \\
    \hline
    RNF-PRF-01 & Reattività UI durante analisi & Performance & MUST \\
    \hline
    RNF-PRF-02 & Bassa latenza in LAN & Performance & SHOULD \\
    \hline
    RNF-MNT-01 & Documentazione Javadoc & Manutenibilità & MUST \\
    \hline
    RNF-MNT-02 & Configurabilità tramite .properties & Manutenibilità & MUST \\
    \hline
    RNF-PRT-01 & Compatibilità JDK 8 & Portabilità & MUST \\
    \hline
    RNF-PRT-02 & Dipendenze auto-contenute & Portabilità & MUST \\
    \hline
    RNF-PRT-03 & Percorsi relativi & Portabilità & MUST \\
    \hline
\end{longtable}

\subsection{Appendice B — Protocollo di Comunicazione (Schema Concettuale)}

Di seguito è riportato uno schema concettuale del protocollo di comunicazione tra client e server (la definizione dettagliata dei messaggi è lasciata al gruppo):

\begin{lstlisting}[language={}, caption={Flusso messaggi client-server}]
CLIENT --> SERVER : LOGIN_REQUEST { username, password }
SERVER --> CLIENT : LOGIN_RESPONSE { success | fail, message }

CLIENT --> SERVER : REGISTER_REQUEST { username, password }
SERVER --> CLIENT : REGISTER_RESPONSE { success | fail, message }

-- Quando entrambi i client sono autenticati --

SERVER --> CLIENT : CHALLENGE_START {
    text_excerpt,
    encrypted_words,
    timeout_seconds
}

CLIENT --> SERVER : CHALLENGE_ANSWER { word_guess }

SERVER --> CLIENT : CHALLENGE_RESULT {
    winner_username,
    correct_word,
    outcome (win | loss | draw)
}
\end{lstlisting}

\subsection{Appendice C — Algoritmo del Cifrario di Cesare}

Il cifrario di Cesare sostituisce ogni lettera dell'alfabeto con la lettera che si trova a uno spostamento (shift) $k$ posizioni avanti nell'alfabeto. La formula per la cifratura è:

\[
    E(c) = (c - \text{`A'} + k) \mod 26 + \text{`A'}
\]

Dove $c$ è il carattere originale e $k$ è lo shift casuale scelto dal server ($1 \leq k \leq 25$). La decifratura usa lo shift inverso $26 - k$.

\begin{lstlisting}[language=Java, caption={Esempio di implementazione Java del cifrario di Cesare}]
public static String encrypt(String word, int shift) {
    StringBuilder sb = new StringBuilder();
    for (char c : word.toCharArray()) {
        if (Character.isLetter(c)) {
            char base = Character.isUpperCase(c) ? 'A' : 'a';
            sb.append((char) ((c - base + shift) % 26 + base));
        } else {
            sb.append(c);
        }
    }
    return sb.toString();
}

public static String decrypt(String word, int shift) {
    return encrypt(word, 26 - shift);
}
\end{lstlisting}

\subsection{Appendice D — Checklist di Verifica Pre-Consegna}

\begin{longtable}{|p{0.5cm}|p{12cm}|p{1.5cm}|}
    \hline
    \rowcolor{primaryBlue!20}
    \textbf{N.} & \textbf{Elemento da verificare} & \textbf{OK?} \\
    \hline
    \endhead
    1 & Codice sorgente completo per server e client & $\square$ \\
    \hline
    2 & Javadoc presente per tutte le classi e metodi pubblici & $\square$ \\
    \hline
    3 & Relazione tecnica con architettura, Use Case, Mockup, Diagrammi & $\square$ \\
    \hline
    4 & Presentazione PowerPoint preparata & $\square$ \\
    \hline
    5 & \code{server.jar} incluso nell'archivio & $\square$ \\
    \hline
    6 & \code{client.jar} incluso nell'archivio & $\square$ \\
    \hline
    7 & File \code{server.properties} incluso e correttamente configurato & $\square$ \\
    \hline
    8 & File \code{client.properties} incluso e correttamente configurato & $\square$ \\
    \hline
    9 & Database SQLite pre-inizializzato con schema e account predefiniti & $\square$ \\
    \hline
    10 & File \code{readme.txt} con credenziali degli account predefiniti & $\square$ \\
    \hline
    11 & Tutte le dipendenze incluse nell'archivio & $\square$ \\
    \hline
    12 & Percorsi di file e DB sono relativi (nessun percorso assoluto) & $\square$ \\
    \hline
    13 & JAR compatibile con JDK 8 & $\square$ \\
    \hline
    14 & Avvio con \code{java -jar server.jar} funzionante senza configurazioni extra & $\square$ \\
    \hline
    15 & Avvio di due istanze \code{java -jar client.jar} funzionante & $\square$ \\
    \hline
    16 & Test completo di una sfida tra due client eseguito con successo & $\square$ \\
    \hline
    17 & Archivio denominato correttamente \code{eseguibili.zip} & $\square$ \\
    \hline
    18 & Consegna effettuata almeno 7 giorni prima dell'appello & $\square$ \\
    \hline
\end{longtable}

% ── Fine documento ──────────────────────────────────────────
\vfill
\begin{center}
    \color{darkGray}\rule{0.5\linewidth}{0.4pt}\\
    \small Fine del documento SRS — GuessTheWord\\
    Versione 1.0 — A.A. 2025/2026
\end{center}

\end{document}
